\documentclass[../../main.tex]{subfiles}
\graphicspath{{../../image/}{../../figure/}}
\begin{document}

\section{消失黏性方法}\label{section:ch04-消失黏性方法}

按照\emph{黏性判据}，\eqref{align:ch04-守恒律组}的解 $U$ 是可容的，只要它是带扩散项的守恒律组的解 $U_{\mu}$ 在 $\mu\downarrow 0$ 时的极限：
\begin{align}
\partial_{t}U(x,t)+\sum_{\alpha=1}^{m}\partial_{\alpha}G_{\alpha}(U(x,t))=\mu\sum_{\alpha,\beta=1}^{m}\partial_{\alpha}[B_{\alpha\beta}(U(x,t))\partial_{\beta}U(x,t)],
\label{align:ch04-黏性组}
\end{align}
其中 $B_{\alpha\beta}$ 是定义于 $\mathcal{O}$ 上的 $n\times n$ 矩阵值函数。

这一原理的动机以及“消失黏性”一词都来自连续物理：正如在前几章看到的，绝热条件下热弹性材料的平衡律诱导一阶双曲型方程组。相反，\cref{chapter:ch03-双曲平衡律方程组}中引入的热黏弹性、导热材料的平衡律生成空间变量中二阶的方程组，含有类似于\eqref{align:ch04-黏性组}右端的扩散项。自然界中，每种材料都在一定程度上具有黏性响应并导热。把一种特定材料归类为弹性非导热体，仅仅意味着黏性与热导率可忽略，尽管并非完全不存在。因而，绝热热弹性理论只有作为热黏弹性当黏性与热导率趋于零时的极限情形，才在物理上有意义。消失黏性方法正是以这一前题为基础。


在建立\eqref{align:ch04-黏性组}时，首要任务是为 $\alpha,\beta=1,\cdots,m$ 选取 $n\times n$ 矩阵 $B_{\alpha\beta}(U)$。对物理起源的方程组，这些矩阵的自然选择由物理决定，或至少为物理所启示。例如，非导热热弹性流体应被视为本构方程\cref{chapter:ch03-双曲平衡律方程组}（2.6.16）、（2.6.17）中黏性与热导率趋于零的 Newton 流体。相应地，当\eqref{align:ch04-守恒律组}代表\cref{chapter:ch03-双曲平衡律方程组}（3.3.29）中不导热（无体积力与供热）热弹性流体的质量、动量、动量与能量平衡律时，带扩散项的相应\eqref{align:ch04-黏性组}应取为\footnote{我们以分量形式写出此方程组，记 $\partial/\partial x_{i}$ 为 $\partial_{i}$，并采用求和约定：重复指标在 $1,2,3$ 范围内求和。}

\begin{align}
\begin{cases}
\partial_{t}\rho+\partial_{j}(\rho v_{j})=0\\
\partial_{t}(\rho v_{i})+\partial_{j}(\rho v_{i}v_{j})+\partial_{i}p(\rho,s)=\lambda\partial_{i}\partial_{j}v_{j}+\mu\partial_{j}(\partial_{i}v_{j}+\partial_{j}v_{i})\\
\partial_{t}\left[\rho\varepsilon(\rho,s)+\tfrac{1}{2}\rho|v|^{2}\right]+\partial_{j}\left[\left(\rho\varepsilon(\rho,s)+\tfrac{1}{2}\rho|v|^{2}+p(\rho,s)\right)v_{j}\right]\\
\qquad=\lambda\partial_{i}(v_{i}\partial_{j}v_{j})+\mu\partial_{j}\left[(\partial_{i}v_{j}+\partial_{j}v_{i})v_{i}\right]+\kappa\partial_{i}\partial_{i}\theta.
\end{cases}
\label{align:ch04-NS}
\end{align}

读者应注意，\eqref{align:ch04-NS}含有三个独立物理参数，即体积黏性 $\lambda+\tfrac{2}{3}\mu$、剪切黏性 $\mu$ 与热导率 $\kappa$，它们都可能很小，但量级不同。因而应准备考虑比\eqref{align:ch04-黏性组}更一般的、含多个独立小参数的消失黏性原理表述。然而，这里将不追求这一推广。

物理启示了\eqref{align:ch04-黏性组}在第 3.3 小节所考虑的每一个例子中的自然形式，包括电磁学、磁流体力学等。另一方面，在缺乏物理指引时，或仅为分析与计算上的方便，可尝试在\eqref{align:ch04-守恒律组}右端添加\emph{人工黏性}。例如，可在\eqref{align:ch04-Burgers}上添加人工黏性以导出\emph{带黏性的 Burgers 方程}：
\begin{align}
\partial_{t}u(x,t)+\partial_{x}\left(\tfrac{1}{2}u^{2}(x,t)\right)=\mu\partial_{x}^{2}u(x,t).
\label{align:ch04-Burgers黏性}
\end{align}

显然，人工黏性的选择应使 $B_{\alpha\beta}$ 诱导耗散，从而使\eqref{align:ch04-黏性组}的 Cauchy 问题适定。诱惑是选用使\eqref{align:ch04-黏性组}抛物化的 $B_{\alpha\beta}$ 矩阵：特别地，当 $\alpha\neq\beta$ 时取零矩阵、当 $\alpha=\beta$ 时取单位矩阵，这将把右端化为 $\mu\Delta U$。然而，\eqref{align:ch04-NS}所物理启示的例子表明，把注意力专门集中于抛物情形是不明智的。一般而言，必须处理中间抛物--双曲型方程组；此时建立 Cauchy 问题的适定性可能需要相当大的努力。见\cref{chapter:ch05-熵与经典解稳定性}。


假定已选取一种消失黏性机制，使 Cauchy 问题\eqref{align:ch04-黏性组}、\eqref{align:ch04-Cauchy初值}适定，于是产生一个问题：当 $\mu\downarrow 0$ 时，解族 $\{U_{\mu}\}$ 在何种收敛意义下收敛？这当然是一个严肃的问题：要求非常强的收敛可能对解族 $\{U_{\mu}\}$ 的紧性提出不合理的期望；另一方面，若收敛意义太弱，则 $\lim U_{\mu}$ 是否为\eqref{align:ch04-守恒律组}的解并不清楚。此问题的各方面将在以后主要于第 VI、XV、XVII 章中讨论。

另一个重要任务是比较消失黏性意义下的可容性与指定熵不等式\eqref{align:ch04-熵不等式}意义下的可容性。在连续热力学中，如第 4.5 小节所讨论的，每当\eqref{align:ch04-黏性组}来自与 Clausius--Duhem 不等式相容的真实本构方程，且\eqref{align:ch04-熵不等式}本身就是或源自 Clausius--Duhem 不等式时，由消失黏性得到的\eqref{align:ch04-守恒律组}的解将自动满足\eqref{align:ch04-熵不等式}。例如，作为\eqref{align:ch04-NS}的解在 $(\lambda,\mu,\kappa)\downarrow 0$ 极限而得到的\cref{chapter:ch03-双曲平衡律方程组}（3.3.29）的解，将自动满足不等式\cref{chapter:ch03-双曲平衡律方程组}（3.3.32）。


若 $\eta(U)$ 是\eqref{align:ch04-守恒律组}的熵，相伴有熵流 $Q(U)$，则\eqref{align:ch04-黏性组}的任一（古典）解 $U_{\mu}$ 满足恒等式
\begin{align}
\partial_{t}\eta(U_{\mu})+\sum_{\alpha=1}^{m}\partial_{\alpha}Q_{\alpha}(U_{\mu})&=\mu\sum_{\alpha,\beta=1}^{m}\partial_{\alpha}[\mathrm{D}\eta(U_{\mu})B_{\alpha\beta}(U_{\mu})\partial_{\beta}U_{\mu}]\nonumber\\
&\quad-\mu\sum_{\alpha,\beta=1}^{m}(\partial_{\alpha}U_{\mu})^{\top}\mathrm{D}^{2}\eta(U_{\mu})B_{\alpha\beta}(U_{\mu})\partial_{\beta}U_{\mu}.
\label{align:ch04-黏性熵恒等式}
\end{align}
右端第二项应为耗散项，使得与 $\mathrm{D}^{2}\eta B_{\alpha\beta}$ 相联系的二次型半正定。除此之外，这一项还须负责压倒\eqref{align:ch04-黏性熵恒等式}右端第一项以及\eqref{align:ch04-黏性组}右端。一个充分（但非必要）条件是
\begin{align}
\sum_{\alpha,\beta=1}^{m}\xi_{\alpha}^{\top}\mathrm{D}^{2}\eta(U)B_{\alpha\beta}(U)\xi_{\beta}\geqslant a\sum_{\alpha=1}^{m}\left|\sum_{\beta=1}^{m}B_{\alpha\beta}(U)\xi_{\beta}\right|^{2},
\label{align:ch04-耗散条件}
\end{align}
其中 $a$ 为某正常数，对任一 $U\in\mathcal{O}$ 与所有 $\xi_{\alpha}\in\mathbb{R}^{n}$，$\alpha=1,\cdots,m$。注意，当 $\alpha\neq\beta$ 时 $B_{\alpha\beta}=0$、当 $\alpha=\beta$ 时 $B_{\alpha\beta}=I$，\eqref{align:ch04-耗散条件}化为 $\eta(U)$ 一致凸的命题。

现设初值 $U_{0}$ 与\eqref{align:ch04-黏性组}、\eqref{align:ch04-Cauchy初值}的解 $U_{\mu}$ 当 $|x|\to\infty$ 时足够快地趋于常状态 $\bar{U}$。不失一般性，可设 $\eta(\bar{U})=0$，$\mathrm{D}\eta(\bar{U})=0$；否则可用\eqref{align:ch04-eta规范化}定义的 $\bar{\eta}(U)$ 替换 $\eta(U)$。进一步假定 $\bar{U}$ 实际是 $\eta$ 在 $\mathcal{O}$ 上的最小值。当 $\eta(U)$ 为凸时这当然自动成立。在这些假设下，在 $\mathbb{R}^{m}\times[0,T)$ 上积分\eqref{align:ch04-黏性熵恒等式}得估计
\begin{align}
\mu\int_{0}^{T}\int_{\mathbb{R}^{m}}\sum_{\alpha,\beta=1}^{m}(\partial_{\alpha}U_{\mu})^{\top}\mathrm{D}^{2}\eta(U_{\mu})B_{\alpha\beta}(U_{\mu})\partial_{\beta}U_{\mu}\,\mathrm{d}x\mathrm{d}t\leqslant\int_{\mathbb{R}^{m}}\eta(U_{0}(x))\,\mathrm{d}x.
\label{align:ch04-黏性熵估计}
\end{align}

我们已为证明黏性可容性判据蕴含熵可容性条件奠定了基础。

\begin{theorem}\label{theorem:ch04-黏性极限熵}
在上述关于 $\eta(U)$ 与 $\{U_{\mu}\}$ 的假设下，设序列 $\{U_{\mu_{k}}\}$，$\mu_{k}\to 0$（当 $k\to\infty$），在 $\mathbb{R}^{m}\times[0,T)$ 上几乎处处有界收敛于函数 $U$。则 $U$ 是\eqref{align:ch04-守恒律组}、\eqref{align:ch04-Cauchy初值}在 $\mathbb{R}^{m}\times[0,T)$ 上的弱解，并满足熵可容性条件\eqref{align:ch04-熵积分0}。
\end{theorem}

\begin{proof}
以任一在 $\mathbb{R}^{m}\times[0,T)$ 内紧支集、取值于 $\mathbb{M}^{1\times n}$ 的 Lipschitz 检验函数 $\Phi(x,t)$ 乘\eqref{align:ch04-黏性组}，在 $\mathbb{R}^{m}\times[0,T)$ 上积分。对时间与空间变量分部积分后，得
\begin{align}
\int_{0}^{T}\int_{\mathbb{R}^{m}}\left[\partial_{t}\Phi U_{\mu}+\sum_{\alpha=1}^{m}\partial_{\alpha}\Phi G_{\alpha}(U_{\mu})\right]\,\mathrm{d}x\mathrm{d}t+\int_{\mathbb{R}^{m}}\Phi(x,0)U_{0}(x)\,\mathrm{d}x\nonumber\\
\quad=\mu\int_{0}^{T}\int_{\mathbb{R}^{m}}\sum_{\alpha=1}^{m}\partial_{\alpha}\Phi B_{\alpha\beta}(U_{\mu})\partial_{\beta}U_{\mu}\,\mathrm{d}x\mathrm{d}t.
\label{align:ch04-黏性弱形式}
\end{align}
由\eqref{align:ch04-耗散条件}与\eqref{align:ch04-黏性熵估计}，当 $\mu_{k}\to 0$ 时右端趋于零，因而极限函数 $U$ 满足\eqref{align:ch04-弱解0}。

接着，以任一在 $\mathbb{R}^{m}\times[0,T)$ 内紧支集的非负 Lipschitz 检验函数 $\psi(x,t)$ 乘\eqref{align:ch04-黏性熵恒等式}，在带形域 $\mathbb{R}^{m}\times[0,T)$ 上积分。对空间与时间变量分部积分后，得到恒等式
\begin{align}
\int_{0}^{T}\int_{\mathbb{R}^{m}}\left[\partial_{t}\psi\eta(U_{\mu})+\sum_{\alpha=1}^{m}\partial_{\alpha}\psi Q_{\alpha}(U_{\mu})\right]\,\mathrm{d}x\mathrm{d}t+\int_{\mathbb{R}^{m}}\psi(x,0)\eta(U_{0}(x))\,\mathrm{d}x\nonumber\\
\quad=\mu\int_{0}^{T}\int_{\mathbb{R}^{m}}\partial_{\alpha}\psi\mathrm{D}\eta(U_{\mu})B_{\alpha\beta}(U_{\mu})\partial_{\beta}U_{\mu}\,\mathrm{d}x\mathrm{d}t\nonumber\\
\qquad+\mu\int_{0}^{T}\int_{\mathbb{R}^{m}}\psi\sum_{\alpha,\beta=1}^{m}(\partial_{\alpha}U_{\mu})^{\top}\mathrm{D}^{2}\eta(U_{\mu})B_{\alpha\beta}(U_{\mu})\partial_{\beta}U_{\mu}\,\mathrm{d}x\mathrm{d}t.
\label{align:ch04-黏性熵弱形式}
\end{align}
由\eqref{align:ch04-耗散条件}与\eqref{align:ch04-黏性熵估计}，右端第一项当 $\mu_{k}\to 0$ 时趋于零，而第二项非负。因此极限函数 $U$ 满足不等式\eqref{align:ch04-熵积分0}。证明完成。
\end{proof}

与黏性判据同类的更一般可容性条件，可通过把\eqref{align:ch04-黏性组}替换为下列形式的方程组来表述：
\begin{align}
\partial_{t}U+\sum_{\alpha=1}^{m}\partial_{\alpha}G_{\alpha}(U)=\mu\sum_{\alpha,\beta=1}^{m}\partial_{\alpha}[B_{\alpha\beta}(U)\partial_{\beta}U]+\nu\sum_{\alpha,\beta,\gamma=1}^{m}\partial_{\alpha}[H_{\alpha\beta\gamma}(U)\partial_{\beta}\partial_{\gamma}U],
\label{align:ch04-黏性色散组}
\end{align}
它涉及三阶、有时甚至四阶微分算子以及两个“消失”参数 $\mu$ 与 $\nu$。例如，可代替\eqref{align:ch04-Burgers黏性}取
\begin{align}
\partial_{t}u(x,t)+\partial_{x}\left(\tfrac{1}{2}u^{2}(x,t)\right)=\mu\partial_{x}^{2}u(x,t)+\nu\partial_{x}^{3}u(x,t).
\label{align:ch04-Burgers三阶}
\end{align}

当黏性的耗散效应与某种色散机制（例如毛细作用所诱导）共存时，物理启示了通过\eqref{align:ch04-黏性色散组}实现可容性的途径。相应地，与\eqref{align:ch04-黏性色散组}相联系的可容性条件称为\emph{黏性--毛细判据}。\eqref{align:ch04-守恒律组}的哪些解能通过这一可容性检验，一般不仅依赖于 $B_{\alpha\beta}$ 与 $H_{\alpha\beta\gamma}$ 的选择，还依赖于 $\mu$ 与 $\nu$ 趋于零的相对速率。作为最低要求，\eqref{align:ch04-黏性色散组}必须与热力学第二定律相容，即对应于\cref{theorem:ch04-黏性极限熵}的命题对物理所提供的熵--熵流对成立。

\end{document}
